\section{From divisor eigenvalues to independent curve degrees}
\label{sec:trace}

We now turn algebraic independence of a joint eigenvalue tuple into a
nonzero product of curve correspondences in independent degrees.  The
point requiring care is that a common triangular basis can depend on
the variables $q_1,\ldots,q_r$.  The following lemma accounts for the derivatives
of that basis and applies without a semisimplicity assumption.

For the Hochschild generalized-trace and inner-conjugation viewpoint,
see Khalkhali~\cite[Section~3.5, Examples~3.5.5--3.5.6 and Lemma~3.5.1]{Khalkhali2009}.
We give an elementary proof of the specific differential-form statement
needed here.

\begin{lemma}[An alternating trace form]
\label{lem:alternating-trace}
Let $L$ be a field extension of $\mathbb C$, let
$A_1,\ldots,A_r\in\operatorname{Mat}_m(L)$ commute pairwise, and let
$A_0\in\operatorname{Mat}_m(L)$ commute with each $A_j$, where $r\geq1$.
Use the K\"ahler differential $\dd$ of $L/\mathbb C$, and multiply
matrix-valued differential forms by matrix multiplication and exterior
multiplication.  Then
\begin{equation}
\label{eq:alternating-trace}
 \Theta(A_0;A_1,\ldots,A_r)
 :=\sum_{\pi\in\mathfrak S_r}\operatorname{sgn}(\pi)
 \operatorname{Tr}\bigl(A_0\,\dd A_{\pi(1)}\cdots \dd A_{\pi(r)}\bigr)
\end{equation}
is unchanged by replacing every $A_j$, including $A_0$, with
$GA_jG^{-1}$ for any $G\in\operatorname{GL}_m(L)$, and differentiating
the new entries.
\end{lemma}

\begin{proof}
Fix an arbitrary matrix-valued one-form $K$. For a scalar parameter $t$,
set
\[
 \begin{aligned}
 B_j(t)&=\dd A_j+t[K,A_j],\\
 F(t)&=\sum_{\pi\in\mathfrak S_r}\operatorname{sgn}(\pi)
       \operatorname{Tr}(A_0B_{\pi(1)}(t)\cdots B_{\pi(r)}(t)).
 \end{aligned}
\]
Thus $F(t)$ is a polynomial with values in differential forms.
Differentiating $[A_j,A_l]=0$ and using the Jacobi identity gives
\begin{equation}
\label{eq:commuting-differentials}
 [A_j,B_l(t)]=[A_l,B_j(t)].
\end{equation}
Here each commutator contains a degree-zero matrix, so it is the ordinary
commutator.

We prove that the derivative $F'(t)$ is
zero separately at each factor position.  Fix a position $s$ and a
permutation $\pi$, and abbreviate $C_h=B_{\pi(h)}(t)$ and $j=\pi(s)$.
Write
\[
 P=C_1\cdots C_{s-1},\qquad Q=C_{s+1}\cdots C_r.
\]
The contribution obtained by differentiating position $s$ is
\[
 T_{\pi,s}=\operatorname{Tr}(A_0PKA_jQ)
              -\operatorname{Tr}(A_0PA_jKQ).
\]
To compute it, move $A_j$ rightward through $Q$, cyclically from the
end of the trace to the beginning, through $A_0$, and then rightward
through $P$.  Expanding every commutator gives the exact identity
\begin{align}
 T_{\pi,s}
 &=\sum_{h<s}\operatorname{Tr}\bigl(
    A_0C_1\cdots C_{h-1}[A_j,C_h]
    C_{h+1}\cdots C_{s-1}KQ\bigr) \notag\\
 &\quad+\sum_{h>s}\operatorname{Tr}\bigl(
    A_0PKC_{s+1}\cdots C_{h-1}[A_j,C_h]
    C_{h+1}\cdots C_r\bigr).
 \label{eq:trace-telescoping}
\end{align}
Empty products have value $I$.  There is no term from crossing $A_0$
because $[A_j,A_0]=0$.  There is also no exterior sign in this identity:
the transported factor $A_j$ has degree zero, and no two one-forms
are exchanged.

For each fixed $h\ne s$, pair a permutation $\pi$ with the permutation
obtained by exchanging its labels at positions $s$ and $h$.
All factors in the corresponding summand of
\eqref{eq:trace-telescoping} remain the same except that
$[A_{\pi(s)},B_{\pi(h)}(t)]$ becomes
$[A_{\pi(h)},B_{\pi(s)}(t)]$.
These are equal by \eqref{eq:commuting-differentials}, while the
permutation sign changes.  Thus
\[
 \sum_{\pi\in\mathfrak S_r}\operatorname{sgn}(\pi)T_{\pi,s}=0
 \qquad\text{for every }s.
\]
The entire alternating trace polynomial is therefore constant in $t$.

Finally, set $K=G^{-1}\dd G$.  Direct differentiation yields
\[
 \dd(GA_jG^{-1})=G\bigl(\dd A_j+[K,A_j]\bigr)G^{-1}.
\]
In a product of these forms the degree-zero matrices $G^{-1}G$ cancel,
and the trace is invariant under conjugation.  Constancy between $t=0$
and $t=1$ proves the assertion.  Notice that the argument never
differentiates $A_0$ and does not require it to commute with $\dd A_j$.
\end{proof}

\begin{lemma}[A nonzero trace from one eigenvalue tuple]
\label{lem:spectral-trace}
Let $L$ be an algebraically closed extension of $\mathbb C$ and let
$A_1,\ldots,A_r\in\operatorname{Mat}_m(L)$ commute pairwise.
Suppose a joint diagonal tuple $(\lambda_1,\ldots,\lambda_r)$ in a
common upper triangular basis is algebraically independent over
$\mathbb C$.  There is a matrix $A_0$, polynomial over $L$ in the
$A_j$, for which \eqref{eq:alternating-trace} is nonzero.
\end{lemma}

\begin{proof}
Let $\Lambda$ be the finite set of distinct joint diagonal tuples,
and write $\lambda=(\lambda_1,\ldots,\lambda_r)$ for the chosen tuple.
For each $\mu\in\Lambda\setminus\{\lambda\}$ choose an index $j(\mu)$
with $\lambda_{j(\mu)}\ne\mu_{j(\mu)}$.  The polynomial
\[
 p(T_1,\ldots,T_r)
 =\prod_{\mu\in\Lambda\setminus\{\lambda\}}
   \frac{T_{j(\mu)}-\mu_{j(\mu)}}
        {\lambda_{j(\mu)}-\mu_{j(\mu)}}
\]
has value $1$ at $\lambda$ and $0$ at every other tuple.  Set
$A_0=p(A_1,\ldots,A_r)$.  This matrix commutes with every $A_j$.
We only use its prescribed diagonal entries; it need not be an
idempotent on a generalized eigenspace.

By Lemma~\ref{lem:alternating-trace}, the form can be computed in the
common upper triangular basis.  Derivatives and products remain upper
triangular, and their diagonal entries multiply.  If the chosen tuple
occurs with multiplicity $m_\lambda$, it follows that
\[
 \Theta(A_0;A_1,\ldots,A_r)
 =r!m_\lambda\,\dd\lambda_1\wedge\cdots\wedge \dd\lambda_r.
\]
The alternating sign and the sign from reordering the scalar one-forms
cancel for each permutation.  Algebraic independence in characteristic
zero makes the displayed wedge nonzero; the integer $r!m_\lambda$ is
also nonzero.  This proves the lemma.
\end{proof}

We apply these lemmas over
$L=\overline{\mathbb C(q_1,\ldots,q_r)}$ to the quantum divisor
operators.  Common triangularization is available for any finite
commuting collection over an algebraically closed field.  The
differential extends to $L$, since the algebraic extension is
separable.

\begin{proposition}[Independent curve correspondences]
\label{prop:independent-correspondences}
Suppose that the quantum divisor operators $A_1(q),\ldots,A_r(q)$
have an algebraically independent joint eigenvalue tuple over
$\mathbb C$.  There are linearly independent numerical degree vectors
$\gamma_1,\ldots,\gamma_r$ of positive genus-zero stable maps such that
\begin{equation}
\label{eq:nonzero-independent-product}
 S_{\gamma_1}\cdots S_{\gamma_r}\ne0,
 \qquad
 \sum_{j=1}^r(d_{\gamma_j}-1)\le n.
\end{equation}
Consequently $r(\iota-1)\le n$.
\end{proposition}

\begin{proof}
Choose $A_0$ by Lemma~\ref{lem:spectral-trace}. In the original fixed
cohomology basis, the classical terms of $A_j$ are constant, and hence
\[
 \dd A_j=\sum_{\gamma>0}q^\gamma\gamma_jS_\gamma\,\omega_\gamma,
 \qquad
 \omega_\gamma=\sum_{l=1}^r\gamma_l\frac{\dd q_l}{q_l}.
\]
These sums are finite. Each ordered summand in the expansion of
\eqref{eq:alternating-trace} is a scalar multiple of
\[
 \operatorname{Tr}(A_0S_{\gamma_1}\cdots S_{\gamma_r})\,
 \omega_{\gamma_1}\wedge\cdots\wedge\omega_{\gamma_r}.
\]
Since $\Theta\ne0$, at least one of these ordered summands is nonzero.
Its trace factor and its wedge factor are therefore both nonzero.
The forms $\dd q_l/q_l$ are an $L$-basis of the differentials of
$L/\mathbb C$, so the nonzero wedge gives linear independence of
the degree vectors. The nonzero trace gives the operator product in
\eqref{eq:nonzero-independent-product}, in the order selected by the
summand.

The operator $S_\gamma$ maps $H_k$ into $H_{k+1-d_\gamma}$. Thus this
product lowers the codimension grading by $\sum_j(d_{\gamma_j}-1)$.
The available gradings run only from $0$ through $n$, proving the upper
bound in \eqref{eq:nonzero-independent-product}. Every positive
stable-map degree has anticanonical degree at least $\iota$, so
\[
 r(\iota-1)\le\sum_{j=1}^r(d_{\gamma_j}-1)\le n.
\]
\end{proof}

Proposition~\ref{prop:spectral} supplies the algebraically independent
joint tuple required in Proposition~\ref{prop:independent-correspondences}.
This proves the inequality in Theorem~\ref{thm:main}; the equality case
is considered next.
